\section{Is low volatility an anomaly here?}
\label{sec:anomaly}
The low-volatility effect has been documented for decades, from
\citet{haugen1991efficient} through \citet{ang2006crosssection} and
\citet{blitz2007volatility}, and two accounts compete for it. One is a
leverage-constraint story: investors who want more market exposure than they
can borrow to obtain bid up high-beta stocks instead, leaving low-beta stocks
cheap \citep{frazzini2014betting}. The other is a benchmarking story: managers
judged against a capitalisation-weighted index treat low-beta positions as
risky in the only sense that matters to them, which deters the arbitrage that
would close the gap \citep{baker2011benchmarks}. Both predict that
low-volatility portfolios earn more than their market risk justifies, and they
differ in why.

Neither question is what a replication study is equipped to settle. What it can
do is ask whether the portfolios it built, and the funds they track, show
returns a standard factor model does not already account for.
Table~\ref{tab:alpha} reports intercepts on the \citet{fama2015five} five
factors plus momentum, then the same regressions with the
betting-against-beta factor added. Standard errors are those of
\citet{newey1987simple}, at eight lags, because daily returns on a
semi-annually rebalanced book are serially correlated and ordinary errors would
overstate significance.

\begin{table}[tbp]
\caption{Six-factor alphas before and after the betting-against-beta factor.}
\label{tab:alpha}
\centering\small
\setlength{\tabcolsep}{5pt}
\begin{tabular}{lrrrrrr}
\toprule
Portfolio & $\alpha$ (\%) & $t$ & $\alpha$+BAB (\%) & $t$ & $\beta_{\rm BAB}$ & $t$\\
\midrule
USMV, the fund & -0.17 & -0.10 & -1.96 & -1.45 & 0.309 & 12.2\\
USMV replica, structural & -0.07 & -0.09 & -0.54 & -0.71 & 0.099 & 10.3\\
USMV replica, sample & 0.58 & 0.66 & 0.16 & 0.19 & 0.113 & 8.8\\
USMV replica, Ledoit--Wolf & 0.91 & 0.96 & 0.51 & 0.58 & 0.116 & 9.6\\
USMV replica, PCA & 0.98 & 1.09 & 0.56 & 0.68 & 0.109 & 8.7\\
USMV replica, beta exempt & 0.34 & 0.28 & -0.84 & -0.80 & 0.241 & 13.6\\
\addlinespace
SPLV, the fund & -1.19 & -0.56 & -3.35 & -1.95 & 0.472 & 14.0\\
SPLV replica & -0.48 & -0.24 & -2.50 & -1.51 & 0.448 & 14.0\\
\bottomrule
\end{tabular}
\notetext{Daily excess returns regressed on the Fama--French five factors plus momentum, then re-run with the betting-against-beta factor added. Annualised intercepts, with Newey--West standard errors at 8 lags. The two blocks end on different dates because the factor vintages do: the six-factor window runs 2011-05-06 to 2026-08-31 and the seven-factor window to 2026-07-31, so the columns are not a paired test on an identical sample. No intercept reaches significance at the five per cent level in either block, though the largest, the SPLV fund's own at $t=-1.95$, falls only just short at $p=0.051$ and would clear a ten per cent threshold. Every betting-against-beta loading is significant by a wide margin, which is why the paper reads the pattern as absorption rather than as an absence of exposure.}
\end{table}


No intercept in either block reaches significance at the five per cent level.
The nearest is the heuristic fund's own, at $t=-1.95$ and $p=0.051$ once
betting-against-beta is included, and it is negative. Every
betting-against-beta loading is significant by a wide margin, with $t$ between
8.7 and 14.0 across every portfolio tested, and adding the factor moves each
intercept downward: the heuristic fund's annualised alpha falls from $-1.19\%$
to $-3.35\%$, the optimised fund's from $-0.17\%$ to $-1.96\%$.

Read carefully, this says something narrower than "the low-volatility anomaly
is explained". It says that over 2011 to 2026, in daily US large-cap data,
these particular portfolios earned no return that the five Fama--French factors
plus momentum fail to account for, and that their exposure to a
betting-against-beta portfolio is large and precisely estimated. A loading is
not a mechanism. The result is consistent with the leverage-constraint account
and it does not distinguish that account from the benchmarking one, because
both predict exactly this exposure; separating them requires evidence on
investor behaviour and constraint bindingness that daily return regressions
cannot supply.

Two further boundaries. The regressions end earlier than the portfolio results
because the factor vintages do, at 2026-08-31 for the six-factor block and
2026-07-31 once betting-against-beta is included, so the two blocks are not a
paired test on an identical sample and the table says so. And an intercept that
is not statistically distinguishable from zero is not an intercept shown to be
zero; with fifteen years of daily data the standard errors here are small, but
the finding is an absence of detected alpha, not a demonstration of its
absence.
